\documentclass[11pt,a4paper]{article}

\usepackage[margin=1in]{geometry}
\usepackage{graphicx}
\usepackage{hyperref}
\usepackage{listings}
\usepackage{xcolor}
\usepackage{amsmath}
\usepackage{amssymb}
\usepackage{booktabs}
\usepackage{enumitem}
\usepackage{float}
\usepackage{titlesec}

\hypersetup{
    colorlinks=true,
    linkcolor=blue!70!black,
    urlcolor=blue!70!black,
    citecolor=blue!70!black,
}

\lstset{
    basicstyle=\ttfamily\small,
    backgroundcolor=\color{gray!8},
    frame=single,
    framerule=0.4pt,
    rulecolor=\color{gray!40},
    breaklines=true,
    showstringspaces=false,
    keywordstyle=\color{blue!70!black}\bfseries,
    commentstyle=\color{green!50!black},
    stringstyle=\color{red!60!black},
    columns=flexible,
    xleftmargin=1em,
    framexleftmargin=0.5em,
}

\title{\textbf{ResNet: Post-Crisis Resource Allocation System} \\[0.5em]
\large Technical Documentation}
\author{}
\date{}

\begin{document}
\maketitle
\tableofcontents
\newpage

% ============================================================================
\section{How to Run}
% ============================================================================

\subsection{Prerequisites}

\begin{itemize}
    \item Python 3.10 or higher
    \item Node.js 18 or higher
    \item (Optional) CUDA-capable GPU or Apple Silicon for accelerated inference
    \item (Optional) OpenAI API key for generating pseudo-labels with GPT-4o
\end{itemize}

\subsection{Installation}

Install Python dependencies:
\begin{lstlisting}[language=bash]
pip install -r requirements.txt
\end{lstlisting}

Install frontend dependencies:
\begin{lstlisting}[language=bash]
cd resnet
npm install
\end{lstlisting}

\subsection{Running the Demo Application}

Start the Flask backend API server:
\begin{lstlisting}[language=bash]
python demo_backend.py --sample 500 --model checkpoints/model.pt --port 8000
\end{lstlisting}

In a separate terminal, start the React frontend dev server:
\begin{lstlisting}[language=bash]
cd resnet
npm run dev
\end{lstlisting}

Open \texttt{http://localhost:5173} in a browser. Click \textbf{Run} to begin the analysis pipeline.

\subsubsection{Backend CLI Options}

\begin{table}[H]
\centering
\begin{tabular}{lll}
\toprule
\textbf{Flag} & \textbf{Default} & \textbf{Description} \\
\midrule
\texttt{--sample} & 500 & Number of posts to load \\
\texttt{--model} & \texttt{checkpoints/model.pt} & Path to trained model checkpoint \\
\texttt{--tsv} & \texttt{CrisisMMD\_v2.0/annotations/...} & Path to CrisisMMD TSV file \\
\texttt{--image-dir} & \texttt{hurricane\_maria} & Directory containing crisis images \\
\texttt{--port} & 8000 & Port for the Flask server \\
\bottomrule
\end{tabular}
\end{table}

\subsection{Training a New Model}

The full training pipeline consists of three steps:

\begin{enumerate}
    \item \textbf{Prepare data} from CrisisMMD:
\begin{lstlisting}[language=bash]
python prepare_data.py --crisismmd CrisisMMD_v2.0 \
    --output data/posts.csv --sample 1000
\end{lstlisting}

    \item \textbf{Generate pseudo-labels} using GPT-4o (requires \texttt{OPENAI\_API\_KEY} in \texttt{.env}):
\begin{lstlisting}[language=bash]
python generate_labels.py --data-dir CrisisMMD_v2.0 \
    --csv data/posts.csv --output data/labels.json
\end{lstlisting}

    \item \textbf{Train the classifier}:
\begin{lstlisting}[language=bash]
python train.py --labels data/labels.json \
    --data-dir CrisisMMD_v2.0 --epochs 200 \
    --output checkpoints/model.pt
\end{lstlisting}
\end{enumerate}

Training produces loss curves and MAE plots in \texttt{checkpoints/plots/}.

\subsection{Building the Frontend for Production}

\begin{lstlisting}[language=bash]
cd resnet
npm run build      # Output to dist/
npm run preview    # Preview the production build
\end{lstlisting}

% ============================================================================
\newpage
\section{Key Innovations}
% ============================================================================

\subsection{GPT-4o Pseudo-Labeling Pipeline}

Traditional crisis datasets require expensive manual annotation. This system uses GPT-4o (via OpenAI API) to generate continuous resource need scores from crisis images and captions automatically.

The labeling prompt provides:
\begin{itemize}
    \item A continuous scoring scale (0.00--1.00) with explicit guidelines to use the full range and avoid binary rounding.
    \item Detailed category definitions for all five resource types.
    \item Six worked examples showing realistic multi-category scoring for different crisis scenarios.
\end{itemize}

Each image is base64-encoded and sent alongside its caption. The script supports incremental progress saving---if interrupted, it resumes from the last labeled sample.

\subsection{Multimodal CLIP Feature Fusion}

Rather than using image or text features alone, the system concatenates both modalities from CLIP ViT-B/32 into a single 1024-dimensional feature vector:

\[
\mathbf{f} = [\text{CLIP}_{\text{image}}(\mathbf{x}) \,\|\, \text{CLIP}_{\text{text}}(c)] \in \mathbb{R}^{1024}
\]

where $\mathbf{x}$ is the crisis image and $c$ is the social media caption. Both embeddings are L2-normalized before concatenation. This captures visual damage information (collapsed structures, flooding) alongside textual context (requests for help, descriptions of conditions) that may not be visible in the image.

\subsection{Real-Time Server-Sent Events (SSE) Streaming}

The inference API uses Server-Sent Events to stream per-post predictions to the frontend in real time, rather than waiting for all posts to be processed. This enables:
\begin{itemize}
    \item Live progress tracking with a progress bar.
    \item Incremental marker animation on the map as each post is analyzed.
    \item Dynamic severity recomputation at the cluster level as new data arrives.
\end{itemize}

\subsection{Interactive Timeline Replay}

After analysis, a slider control allows scrubbing through posts chronologically. Moving the slider recomputes cluster-level severity on the fly using only the posts up to the selected point, visualizing how resource demands evolve as a disaster unfolds.

\subsection{Population-Weighted Geographic Simulation}

For the Hurricane Maria demo, posts are assigned coordinates using population-weighted clustering across Puerto Rico's major cities. A ray-casting point-in-polygon algorithm ensures all generated coordinates fall within the island's coastline, producing realistic geographic distributions.

% ============================================================================
\newpage
\section{System Architecture}
% ============================================================================

\subsection{Overview}

The system consists of four major components:

\begin{enumerate}
    \item \textbf{Data Preparation}: CrisisMMD dataset parsing, GPT-4o pseudo-labeling.
    \item \textbf{Training Pipeline}: CLIP feature encoding, ResourceClassifier training.
    \item \textbf{Backend API}: Flask REST server with SSE streaming for real-time inference.
    \item \textbf{Frontend}: React + TypeScript interactive map with Leaflet.
\end{enumerate}

\subsection{Data Flow}

\begin{verbatim}
CrisisMMD TSV annotations
        |  prepare_data.py
        v
CSV (image_path, caption, timestamp)
        |  generate_labels.py (GPT-4o)
        v
JSON labels (image_path, caption, {infrastructure: 0.XX, ...})
        |  train.py (CLIP encoding + MLP training)
        v
Trained model checkpoint (model.pt)
        |  demo_backend.py
        v
Flask API (REST + SSE) --> React Frontend (Leaflet map)
\end{verbatim}

% ============================================================================
\newpage
\section{Model Architecture}
% ============================================================================

\subsection{CLIP Encoder}

The system uses OpenAI's CLIP ViT-B/32 model (loaded via HuggingFace Transformers) as a frozen feature extractor. No CLIP parameters are fine-tuned during training.

\begin{itemize}
    \item \textbf{Image encoding}: Input image $\rightarrow$ ViT-B/32 $\rightarrow$ 512-dim L2-normalized embedding.
    \item \textbf{Text encoding}: Input caption $\rightarrow$ CLIP text encoder $\rightarrow$ 512-dim L2-normalized embedding (truncated to 77 tokens).
    \item \textbf{Feature fusion}: Concatenation $\rightarrow$ 1024-dim feature vector.
\end{itemize}

\subsection{ResourceClassifier}

A multi-layer perceptron (MLP) that maps 1024-dim CLIP features to five resource need scores:

\begin{table}[H]
\centering
\begin{tabular}{lll}
\toprule
\textbf{Layer} & \textbf{Dimensions} & \textbf{Activation} \\
\midrule
Linear + ReLU + Dropout(0.3) & $1024 \rightarrow 512$ & ReLU \\
Linear + ReLU + Dropout(0.3) & $512 \rightarrow 256$ & ReLU \\
Linear + Sigmoid & $256 \rightarrow 5$ & Sigmoid \\
\bottomrule
\end{tabular}
\end{table}

Output categories (each in $[0, 1]$):
\begin{enumerate}
    \item \textbf{Infrastructure}: Roads, bridges, buildings, power lines, utilities.
    \item \textbf{Food}: Hunger, food shortages, supply disruption.
    \item \textbf{Shelter}: Displacement, destroyed housing, homelessness.
    \item \textbf{Sanitation \& Water}: Contaminated water, flooding, sewage.
    \item \textbf{Medication}: Injuries, hospital damage, medical supply needs.
\end{enumerate}

\subsection{Training Configuration}

\begin{table}[H]
\centering
\begin{tabular}{ll}
\toprule
\textbf{Parameter} & \textbf{Value} \\
\midrule
Loss function & Binary Cross-Entropy (BCE) \\
Optimizer & Adam ($\text{lr} = 10^{-3}$, weight decay $= 10^{-3}$) \\
LR scheduler & ReduceLROnPlateau (factor $= 0.5$, patience $= 7$) \\
Early stopping & Patience $= 20$ epochs on validation loss \\
Batch size & 32 \\
Max epochs & 200 \\
Train/Val split & 80\%/20\% \\
Evaluation metric & Per-category MAE + overall MAE \\
Feature caching & CLIP embeddings cached to \texttt{cached\_features.pt} \\
\bottomrule
\end{tabular}
\end{table}

The training script generates four plots:
\begin{itemize}
    \item Training vs.\ validation loss curve
    \item Overall validation MAE curve
    \item Per-category validation MAE curves
    \item Final per-category MAE bar chart
\end{itemize}

% ============================================================================
\newpage
\section{Spatial Clustering}
% ============================================================================

\subsection{DBSCAN Algorithm}

Geographic clustering uses DBSCAN (Density-Based Spatial Clustering of Applications with Noise) on latitude/longitude coordinates:

\begin{itemize}
    \item \textbf{Training/inference pipeline} (\texttt{clustering.py}): $\varepsilon = 0.01$ (approximately 1\,km), \texttt{min\_samples} $= 2$.
    \item \textbf{Demo backend} (\texttt{demo\_backend.py}): $\varepsilon = 0.15$ (larger radius for the Hurricane Maria demo), \texttt{min\_samples} $= 3$.
\end{itemize}

Noise points (DBSCAN label $= -1$) are discarded from cluster aggregation but still displayed on the map.

\subsection{Cluster-Level Aggregation}

After per-post inference, cluster-level resource demands are computed as the arithmetic mean of all member posts' scores for each category. Severity is determined by a weighted average:

\[
\text{weighted\_severity} = \frac{1}{|C|} \sum_{p \in C} w(s_p), \quad
w(s) = \begin{cases}
0.1 & s = \text{little\_or\_none} \\
0.3 & s = \text{mild} \\
1.0 & s = \text{severe}
\end{cases}
\]

The combined severity label is then assigned:
\begin{itemize}
    \item $< 0.20$: \texttt{little\_or\_none}
    \item $0.20$--$0.30$: \texttt{mild}
    \item $> 0.30$: \texttt{severe}
\end{itemize}

% ============================================================================
\newpage
\section{Backend API}
% ============================================================================

\subsection{Endpoints}

The Flask backend (\texttt{demo\_backend.py}) exposes the following endpoints:

\begin{table}[H]
\centering
\begin{tabular}{lll}
\toprule
\textbf{Method} & \textbf{Path} & \textbf{Description} \\
\midrule
GET & \texttt{/} & Serve frontend HTML \\
GET & \texttt{/api/posts} & Return all posts with DBSCAN cluster labels \\
GET & \texttt{/api/predict} & SSE stream: CLIP encoding + model inference \\
GET & \texttt{/images/<path>} & Serve crisis images from the image directory \\
\bottomrule
\end{tabular}
\end{table}

\subsection{\texttt{GET /api/posts} Response}

Returns a JSON object containing:
\begin{itemize}
    \item \texttt{posts}: Array of post objects, each with \texttt{lat}, \texttt{lon}, \texttt{caption}, \texttt{timestamp}, \texttt{date}, \texttt{image} (URL), and \texttt{cluster} label.
    \item \texttt{categories}: List of the five resource category names.
    \item \texttt{clusters}: Map from cluster ID to metadata (name of nearest city, centroid coordinates, member count).
\end{itemize}

\subsection{\texttt{GET /api/predict} SSE Stream}

Streams two types of events:

\begin{enumerate}
    \item \textbf{Progress events}: Sent after each post is processed.
\begin{lstlisting}
data: {"type": "progress", "current": 42, "total": 500,
       "cluster": 3, "severity_label": "mild",
       "scores": {"infrastructure": 0.72, ...}}
\end{lstlisting}

    \item \textbf{Done event}: Sent after all posts are processed.
\begin{lstlisting}
data: {"type": "done",
       "cluster_scores": {"0": {"infrastructure": 0.75,
            "weighted_severity": 0.403,
            "combined_severity": "mild", ...}, ...}}
\end{lstlisting}
\end{enumerate}

\subsection{Coordinate Assignment}

For the Hurricane Maria demo, posts are assigned synthetic coordinates based on population-weighted cluster centers across eight Puerto Rico cities. A hash of each tweet's ID deterministically maps it to a city, with Gaussian jitter constrained to fall within a simplified polygon of Puerto Rico's coastline (30 vertices, ray-casting point-in-polygon test).

% ============================================================================
\newpage
\section{Frontend}
% ============================================================================

\subsection{Technology Stack}

\begin{table}[H]
\centering
\begin{tabular}{ll}
\toprule
\textbf{Library} & \textbf{Purpose} \\
\midrule
React 19 & UI framework \\
TypeScript 5.9 & Type safety \\
Vite 7 & Build tool and dev server \\
Leaflet + react-leaflet & Interactive map rendering \\
Zustand & Lightweight state management \\
Tailwind CSS 4 & Utility-first styling \\
\bottomrule
\end{tabular}
\end{table}

\subsection{User Interface}

The frontend displays a dark-themed Leaflet map (ArcGIS World Dark Gray basemap) centered on Puerto Rico. Key UI elements:

\begin{itemize}
    \item \textbf{Post markers}: Color-coded circles by severity---green (little/none), yellow (mild), red (severe)---with CSS glow effects for visual emphasis.
    \item \textbf{Control panel}: Bottom-center overlay with a three-step phase indicator (Collect + Cluster $\rightarrow$ CLIP + Model $\rightarrow$ Results), progress bar, status text, and run button.
    \item \textbf{Timeline slider}: Appears after analysis; scrubs through posts chronologically with on-the-fly severity recomputation.
    \item \textbf{Cluster popup modal}: Clicking a cluster marker opens a detail panel showing averaged resource demands and individual posts with their scores, sortable by severity or category.
\end{itemize}

\subsection{Real-Time Animation}

Posts appear on the map via a drip-feed animation system:
\begin{itemize}
    \item An interval timer adds \texttt{POSTS\_PER\_TICK} (5) posts every \texttt{INTERVAL\_MS} (20\,ms).
    \item The drip target advances as SSE progress events arrive from the backend.
    \item Cluster severity colors update dynamically as new posts are assigned.
\end{itemize}

\subsection{Vite Proxy Configuration}

The Vite dev server proxies \texttt{/api/*} and \texttt{/images/*} requests to the Flask backend on port 8000, with SSE buffering disabled to ensure real-time event streaming.

% ============================================================================
\newpage
\section{Data Pipeline}
% ============================================================================

\subsection{CrisisMMD Dataset}

The system uses the CrisisMMD v2.0 dataset, a multimodal crisis dataset containing social media posts from seven major 2017 disasters:

\begin{enumerate}
    \item Hurricane Harvey
    \item Hurricane Irma
    \item Hurricane Maria
    \item Mexico Earthquake
    \item Iraq-Iran Earthquake
    \item California Wildfires
    \item Sri Lanka Floods
\end{enumerate}

Each TSV annotation file contains columns including \texttt{tweet\_id}, \texttt{image\_path}, \texttt{tweet\_text}, \texttt{humanitarian\_category}, and \texttt{damage\_severity}.

\subsection{\texttt{prepare\_data.py}}

Converts CrisisMMD TSV annotations into a simplified CSV format with columns: \texttt{image\_path}, \texttt{caption}, \texttt{timestamp}. Timestamps are extracted from Twitter snowflake IDs. Supports random sampling with a fixed seed for reproducibility.

\subsection{\texttt{generate\_labels.py}}

For each image--caption pair, sends a base64-encoded image and structured prompt to GPT-4o-mini, receiving back a JSON object with five continuous scores. Key design decisions:
\begin{itemize}
    \item Scores are continuous in $[0, 1]$ rather than binary, enabling regression training.
    \item The prompt explicitly discourages binary rounding (0 or 1) to encourage nuanced scoring.
    \item Progress is saved after every sample for fault tolerance.
    \item A configurable delay between API calls prevents rate limiting.
\end{itemize}

\subsection{Feature Caching}

During training, CLIP encodings are computed once and cached to \texttt{data/cached\_features.pt}. Subsequent training runs load features from cache, skipping the expensive CLIP forward pass.

% ============================================================================
\newpage
\section{Dependencies}
% ============================================================================

\subsection{Python}

\begin{table}[H]
\centering
\begin{tabular}{ll}
\toprule
\textbf{Package} & \textbf{Purpose} \\
\midrule
\texttt{torch}, \texttt{torchvision} & Neural network training and inference \\
\texttt{transformers} & CLIP ViT-B/32 model loading \\
\texttt{Pillow} & Image loading and preprocessing \\
\texttt{scikit-learn} & DBSCAN clustering \\
\texttt{numpy} & Numerical operations \\
\texttt{openai} & GPT-4o API for pseudo-labeling \\
\texttt{python-dotenv} & Environment variable loading \\
\texttt{matplotlib} & Training metric plots \\
\texttt{flask} & REST API server \\
\texttt{flask-cors} & Cross-origin request support \\
\bottomrule
\end{tabular}
\end{table}

\subsection{Node.js / Frontend}

\begin{table}[H]
\centering
\begin{tabular}{ll}
\toprule
\textbf{Package} & \textbf{Purpose} \\
\midrule
\texttt{react}, \texttt{react-dom} (19.x) & UI framework \\
\texttt{react-leaflet}, \texttt{leaflet} & Map rendering \\
\texttt{zustand} (5.x) & State management \\
\texttt{@tanstack/react-query} (5.x) & Data fetching \\
\texttt{react-router-dom} (7.x) & Client-side routing \\
\texttt{vite} (7.x) & Build tool \\
\texttt{typescript} (5.9) & Type checking \\
\texttt{tailwindcss} (4.x) & Utility CSS \\
\bottomrule
\end{tabular}
\end{table}

% ============================================================================
\newpage
\section{File Reference}
% ============================================================================

\begin{table}[H]
\centering
\begin{tabular}{ll}
\toprule
\textbf{File} & \textbf{Purpose} \\
\midrule
\texttt{model.py} & ResourceClassifier MLP definition \\
\texttt{encoder.py} & CLIPEncoder wrapper (image + text + cluster encoding) \\
\texttt{train.py} & Training loop with early stopping, LR scheduling, plots \\
\texttt{clustering.py} & DBSCAN spatial clustering with Cluster dataclass \\
\texttt{data\_collector.py} & Post and Call dataclass definitions \\
\texttt{prepare\_data.py} & CrisisMMD TSV $\rightarrow$ CSV converter \\
\texttt{generate\_labels.py} & GPT-4o pseudo-label generator \\
\texttt{demo\_backend.py} & Flask API server with SSE streaming \\
\texttt{main.py} & Standalone inference demo script \\
\texttt{requirements.txt} & Python dependency list \\
\midrule
\texttt{resnet/src/components/map/Map.tsx} & Main map view with inference orchestration \\
\texttt{resnet/src/components/map/cluster.tsx} & Post markers with severity glow effects \\
\texttt{resnet/src/components/map/ClusterPopup.tsx} & Cluster detail modal \\
\texttt{resnet/src/types/post.ts} & Post and AnalyzedPost TypeScript interfaces \\
\texttt{resnet/src/types/cluster.ts} & Cluster TypeScript interface \\
\texttt{resnet/vite.config.ts} & Vite dev server with API proxy configuration \\
\bottomrule
\end{tabular}
\end{table}

\end{document}
